% !TeX root = ../大学生数学竞赛概率论与数理统计试题选编与专题训练.tex
% 由主文档通过 \input 载入；请编译主文档。
\setcounter{chapter}{0} % 保留原书章号。
\chapter{全国大学生数学竞赛概率相关真题}

本章按届次列出十三道概率相关真题。题面与基础解答由 Jacobi Lie 的公开题解转写并校订，属于公开题解转录层级；原卷题号由其逐题元数据给出，未声称逐份校验原卷扫描件。原作者的解答、评注及方法稿采用 CC BY-SA 4.0，具体署名、固定版本及许可见全书编写说明。

\section{第5届决赛：9题}\label{pr:fin05-hi:9}

\begin{example}[公开题解转录]\label{q:fin05-hi:9}

盒中有编号 $1,\ldots,N$ 的 $N$ 张票，每次等概率抽取一张后放回，各次独立。求收集到 $r$ 种不同号码所需抽取次数的期望，其中 $1\le r\le N$。

\end{example}

\href{https://github.com/lzxcr/cmc-mathematics/blob/64916fa1fdbc3e444a652eca204d5669058ee321/src/problems/finals/05/10.md}{逐题来源及版本}。

\begin{solve}

已有 $k$ 种号码时，下一次抽到新号码的概率为 $(N-k)/N$。等待下一个新号码所需次数服从成功概率为 $(N-k)/N$ 的几何分布，其期望为 $N/(N-k)$。由期望的可加性，所求为
\[
\mathbb E T_r=\sum_{k=0}^{r-1}\frac{N}{N-k}
=N(H_N-H_{N-r}),
\]
其中 $H_m=\sum_{j=1}^m1/j$，约定 $H_0=0$。

\end{solve}

\begin{remark}[方法与条件]

这是集券问题。全部收齐的期望为 $NH_N\sim N\log N$；若只要求固定比例 $r/N\to\theta\in(0,1)$，则 $\mathbb E T_r/N\to-\log(1-\theta)$。最后少量缺失号码造成了额外的对数代价。

\end{remark}

\section{第6届决赛：14题}\label{pr:fin06-hi:14}

\begin{example}[公开题解转录]\label{q:fin06-hi:14}

设在国际市场上对我国某种出口商品每年的需求量 $X$（单位：吨）是随机变量，$X$ 服从 $[100,200]$ 上的均匀分布。每出售这种商品一吨，可以为国家挣得外汇 3 万元；若销售不出而囤积于仓库，则每吨需要花费保养费用 1 万元。问应组织多少吨货源，才能使国家的收益最大？

\end{example}

\href{https://github.com/lzxcr/cmc-mathematics/blob/64916fa1fdbc3e444a652eca204d5669058ee321/src/problems/finals/06/15.md}{逐题来源及版本}。

\begin{solve}

以期望收益为目标。备货 $t\ge0$ 吨时，收益可统一写成
\[
Y_t=3t-4(t-X)_+,
\qquad (u)_+=\max(u,0).
\]
记 $F$ 为需求量的分布函数。由于 $X$ 连续分布，期望收益 $R(t)=\mathbb EY_t$ 可微，且
\[
R'(t)=3-4F(t).
\]
这个导数随 $t$ 不增加；在 $t<100$ 时为 $3$，在 $t>200$ 时为 $-1$。因此唯一最优点满足
\[
F(t)=\frac34,
\qquad \frac{t-100}{100}=\frac34,
\]
即应组织 $\boxed{175}$ 吨货源。相应最大期望收益为
\[
R(175)=3\cdot175-\frac4{100}\int_{100}^{175}(175-x)\,\mathrm dx
=412.5\text{ 万元}.
\]

\end{solve}

\begin{remark}[方法与条件]

这是报童模型的分位数原则。若售出一单位收益 $r>0$，过剩一单位费用 $c>0$，则
\[
R(t)=rt-(r+c)\mathbb E(t-X)_+.
\]
连续需求下，内点最优量满足 $F(t)=r/(r+c)$。允许需求有原子时，导数改为左右导数，准确条件是
\[
F(t-)\le\frac r{r+c}\le F(t).
\]
所以最优备货量由边际收益与边际损失决定，是需求的一个分位数，通常并非平均需求。

\end{remark}

\section{第7届决赛：13题}\label{pr:fin07-hi:13}

\begin{example}[公开题解转录]\label{q:fin07-hi:13}

甲袋中有 $N-1$（$N>1$）个白球和 1 个黑球，乙袋中有 $N$ 个白球，每次从甲乙两袋中分别取出一个球并交换放入另一袋中，这样经过了 $n$ 次，求黑球出现在甲袋中的概率 $p_{n}$，并计算 $\lim_{n\to\infty}p_{n}$。

\end{example}

\href{https://github.com/lzxcr/cmc-mathematics/blob/64916fa1fdbc3e444a652eca204d5669058ee321/src/problems/finals/07/15.md}{逐题来源及版本}。

\begin{solve}

黑球每次以概率 $1/N$ 换袋。由全概率公式，
\[
p_{n+1}=\left(1-\frac1N\right)p_n+\frac1N(1-p_n),\qquad p_0=1.
\]
减去平衡值 $1/2$，得到等比递推，故
\[
p_n=\frac12+\frac12\left(1-\frac2N\right)^n,\qquad \lim_{n\to\infty}p_n=\frac12.
\]
这里 $N\geq2$；当 $N=2$ 时，对所有 $n\geq1$ 均有 $p_n=1/2$。

\end{solve}

\begin{remark}[方法与条件]

减去 $1/2$ 的原因是它是两态转移的平衡概率；剩余偏差每步乘以 $1-2/N$，直接给出趋近平衡的速度。这也是线性递推中先找不动点、再研究偏差的最简单例子，见方法卷：离散迭代。

\end{remark}

\section{第8届决赛：13题}\label{pr:fin08-hi:13}

\begin{example}[公开题解转录]\label{q:fin08-hi:13}

（概率统计）设 $\left\{X_{n}\right\}$ 是独立同分布的随机变量序列，且 $P\left(X_{n}=0\right)=P\left(X_{n}=a\right)=\frac{1}{2}$，其中常数 $a>0$。记 $Y_{n}=\sum_{k=1}^{n}\frac{X_{k}}{2^{k}}$，求 $Y_{n}$ 的特征函数，并证明其分布收敛于区间 $[0,a]$ 上的均匀分布。

\end{example}

\href{https://github.com/lzxcr/cmc-mathematics/blob/64916fa1fdbc3e444a652eca204d5669058ee321/src/problems/finals/08/15.md}{逐题来源及版本}。

\begin{solve}

独立公平的二进制位使 $Y_n$ 等概率地取值于 $\{aj/2^n:0\leq j<2^n\}$。所以其特征函数为
\[
\varphi_n(t)=2^{-n}\sum_{j=0}^{2^n-1}e^{iatj/2^n}
=e^{iat(1-2^{-n})/2}\frac{\sin(at/2)}{2^n\sin(at/2^{n+1})},
\]
右端在分母为零处按连续延拓理解。前一个表达式是连续函数 $x\mapsto e^{iatx}$ 在 $[0,1]$ 上的 Riemann 和，故
\[
\varphi_n(t)\longrightarrow\int_0^1e^{iatx}\,\mathrm dx
=\begin{cases}(e^{iat}-1)/(iat),&t\ne0,\\1,&t=0.\end{cases}
\]
极限是 $[0,a]$ 上均匀分布的特征函数，且在零点连续；由 Lévy 连续性定理得到所求分布收敛。

\end{solve}

\begin{remark}[方法与条件]

这里也可直接检验弱收敛：对每个连续函数 $h$，$Eh(Y_n)=2^{-n}\sum_{j<2^n}h(aj/2^n)\to a^{-1}\int_0^ah(u)\,\mathrm du$。换成 $b$ 进制时，独立均匀数字同样产生 $[0,a]$ 上的均匀分布。特征函数与分布的关系见\hyperref[app:probability-transforms]{方法卷：概率变换}。

\end{remark}

\section{第9届决赛：13题}\label{pr:fin09-hi:13}

\begin{example}[公开题解转录]\label{q:fin09-hi:13}

（概率统计）设 $X_1,X_2,\cdots,X_n$ 是独立同分布的随机变量，其有共同的分布函数 $F(x)$ 和密度函数 $f(x)$。现对随机变量 $X_1,X_2,\cdots,X_n$ 按大小顺序重新排列为 $X_{n1} \leq X_{n2} \leq \cdots \leq X_{nn}$。

（1）求随机变量 $(X_{n1},X_{nn})$ 的联合概率密度函数 $f_{1n}(x,y)$；

（2）如果 $X_i\ (i = 1,2,\cdots,n)$ 服从区间 $[0,1]$ 上的均匀分布，求随机变量 $U = X_{nn} + X_{n1}$ 的密度函数 $f_U(u)$。

\end{example}

\href{https://github.com/lzxcr/cmc-mathematics/blob/64916fa1fdbc3e444a652eca204d5669058ee321/src/problems/finals/09/15.md}{逐题来源及版本}。

\begin{solve}

（1）先设 $n\geq2$。最小值位于 $\mathrm dx$、最大值位于 $\mathrm dy$，有 $n(n-1)$ 种指定这两个样本的方式，其余样本须落在 $(x,y)$。于是
\[
f_{1n}(x,y)=\begin{cases}n(n-1)[F(y)-F(x)]^{n-2}f(x)f(y),&x<y,\\0,&x\geq y.\end{cases}
\]
该式也可由 $F(y)^n-[F(y)-F(x)]^n$ 对 $x,y$ 求混合导数得到。

（2）所求变量是两端点之和，而原密度只依赖于间距，因此取 $u=x+y,v=y-x$。逆变换 $x=(u-v)/2,y=(u+v)/2$ 的 Jacobian 绝对值为 $1/2$；由 $0<x<y<1$，积分区域为 $0<v<\min(u,2-u)$。因此
\[
f_U(u)=\begin{cases}\frac n2u^{n-1},&0<u<1,\\\frac n2(2-u)^{n-1},&1<u<2,\\0,&\text{其他}.\end{cases}
\]
若 $n=1$，$(X_{11},X_{11})$ 集中在对角线上，没有二维密度；第二问则为 $U=2X_1$，在 $(0,2)$ 上密度为 $1/2$。

\end{solve}

\begin{remark}[方法与条件]

最小值与最大值一般不独立。联合密度中的 $[F(y)-F(x)]^{n-2}$ 恰记录其余样本全部夹在两端点之间的概率，因此先求联合分布，再变换到“中心与跨度”最为直接。见\hyperref[app:probability-law]{方法卷：分布与变量变换}。

\end{remark}

\section{第10届决赛：13题}\label{pr:fin10-hi:13}

\begin{example}[公开题解转录]\label{q:fin10-hi:13}

（满分 10 分）（概率统计）设独立随机变量序列 $\{X_n,n\geq1\}$ 满足 $P(X_n=\pm n^{\theta})=\frac{1}{2}$，其中 $\theta>0$ 是常数。记 $S_n=\sum_{i=1}^{n}X_i$。

（1）当 $\theta<\frac{1}{2}$ 时，证明 $\frac{S_n}{n}$ 依概率收敛于 $0$，即对任意 $\epsilon>0$，$\lim_{n\to\infty}P\left(\frac{\left|S_n\right|}{n}\geq\epsilon\right)=0$；

（2）证明：$\frac{S_n}{\sqrt{\operatorname{Var}\left(S_n\right)}}\xrightarrow{D}N(0,1)$，其中 $\operatorname{Var}\left(S_n\right)$ 表示 $S_n$ 的方差，$\xrightarrow{D}$ 表示依分布收敛。

\end{example}

\href{https://github.com/lzxcr/cmc-mathematics/blob/64916fa1fdbc3e444a652eca204d5669058ee321/src/problems/finals/10/15.md}{逐题来源及版本}。

\begin{solve}

记 $s_n^2=\operatorname{Var}S_n=\sum_{k=1}^nk^{2\theta}$，由 Riemann 和，
\[
s_n^2\sim\frac{n^{2\theta+1}}{2\theta+1}.
\]
（1）当 $\theta<1/2$ 时，$s_n^2/n^2\to0$。Chebyshev 不等式给出 $P(|S_n|\geq n\varepsilon)\leq s_n^2/(n^2\varepsilon^2)\to0$。

（2）对任意 $\theta>0$，
\[
\max_{k\leq n}\frac{|X_k|}{s_n}=\frac{n^\theta}{s_n}\longrightarrow0.
\]
这说明没有单个随机增量能占据总标准差的固定比例。因此对每个固定 $\varepsilon>0$，Lindeberg 条件中的所有截断项 $E[X_k^2\mathbf1_{\{|X_k|>\varepsilon s_n\}}]$ 在充分大的 $n$ 时均为零。由 Lindeberg 中心极限定理，$S_n/s_n\Rightarrow N(0,1)$。

\end{solve}

\begin{remark}[方法与条件]

中心极限定理还给出精确阈值：$\theta=1/2$ 时 $S_n/n\Rightarrow N(0,1/2)$；$\theta>1/2$ 时 $n/s_n\to0$，故 $P(|S_n|\leq n\varepsilon)\to0$，因为标准正态在零点无原子。归一化尺度与 Lindeberg 条件见\hyperref[app:probability-limit]{方法卷：概率极限定理}。

\end{remark}

\section{第11届决赛：12题}\label{pr:fin11-high:12}

\begin{example}[公开题解转录]\label{q:fin11-high:12}

设 $\left\{X_{i}\right\}$ 是独立随机变量序列。

（1）若 $\left\{X_{i}\right\}$ 服从大数定律且满足中心极限定理，即
\[
\frac{1}{n}\sum_{i=1}^{n}\left(X_{i}-EX_{i}\right)\overset{P}{\rightarrow}0 \text{ 和 } \frac{\sum_{i=1}^{n}\left(X_{i}-EX_{i}\right)}{\sqrt{\sum_{k=1}^{n}\operatorname{Var}\left(X_{k}\right)}}\overset{D}{\rightarrow}N(0,1),
\]
则 $\lim_{n\to\infty}\frac{1}{n^{2}}\sum_{k=1}^{n}\operatorname{Var}\left(X_{k}\right)=0$。

（2）若 $\left\{X_{i}\right\}$ 同分布且满足 $\lim_{n\to\infty}nP\left(\left|X_{1}\right|\geq n\right)=0$，则 $\frac{1}{n}\sum_{i=1}^{n}X_{i}-\mu_{n}$ 依概率收敛于 0，即 $\frac{1}{n}\sum_{i=1}^{n}X_{i}-\mu_{n}\overset{P}{\rightarrow}0$，其中 $\mu_{n}=E\left[X_{1}I\left(\left|X_{1}\right|\leq n\right)\right]$，$I(A)$ 表示事件 $A$ 的示性函数。

\end{example}

\href{https://github.com/lzxcr/cmc-mathematics/blob/64916fa1fdbc3e444a652eca204d5669058ee321/src/problems/finals/11/14.md}{逐题来源及版本}。

\begin{solve}

（1）记 $s_n^2=\sum_{k=1}^n\operatorname{Var}X_k$、$Z_n=\sum_{k=1}^n(X_k-EX_k)/s_n$。若 $s_n/n$ 不趋于零，则某子列上 $s_n/n\geq c>0$，从而
\[
P\left(\left|\frac1n\sum_{k=1}^n(X_k-EX_k)\right|>c\right)\geq P(|Z_n|>1)\longrightarrow P(|Z|>1)>0,
\]
与大数定律矛盾，所以 $s_n^2/n^2\to0$。

（2）尾条件控制“至少一个样本很大”的概率，截断后的方差则可用 Chebyshev 控制；两步必须同时满足。令 $Y_{ni}=X_i\mathbf1_{\{|X_i|\leq n\}}$，则 $EY_{ni}=\mu_n$。由并集估计，$P(\exists i\leq n:X_i\ne Y_{ni})\leq nP(|X_1|>n)\to0$。记 $T=|X_1|$，由题设及单调性，$tP(T>t)\to0$，所以
\[
\frac{EY_{n1}^2}{n}\leq\frac2n\int_0^ntP(T>t)\,\mathrm dt\longrightarrow0.
\]
因此 Chebyshev 不等式给出
\[
P\left(\left|\frac1n\sum_{i=1}^n(Y_{ni}-\mu_n)\right|>\varepsilon\right)
\leq\frac{EY_{n1}^2}{n\varepsilon^2}\longrightarrow0.
\]
结合截断前后不同的概率趋于零，得到所求结论。

\end{solve}

\begin{remark}[方法与条件]

这里不要求 $E|X_1|<\infty$，所以中心必须选截断均值 $\mu_n$，不能未经验证替换成 $EX_1$。尾部截断将大数问题分成稀有大值与可控方差两部分，见\hyperref[app:probability-limit]{方法卷：概率极限定理}。

\end{remark}

\section{第12届决赛：12题}\label{pr:fin12-high:12}

\begin{example}[公开题解转录]\label{q:fin12-high:12}

设有一列盒子，已知第 $k$ 个盒子中有 $k$ 个球，其中 1 个是红球，另外 $k-1$ 个是白球。现相互独立地从前 $n$ 个盒子中各取一球，记 $S_n$ 表示取出的 $n$ 个球中红球的个数。证明：

（1）$\frac{S_n}{\ln(n)}$ 依概率收敛于 1；

（2）$\frac{S_n - \ln(n)}{\sqrt{\ln(n)}}$ 依分布收敛于标准正态分布 $N(0,1)$；

（3）对任意 $r > 0$，$\lim_{n \to \infty} E\left(\frac{\left|S_n - \ln(n)\right|^r}{\ln^r(n) + \left|S_n - \ln(n)\right|^r}\right) = 0$。

\end{example}

\href{https://github.com/lzxcr/cmc-mathematics/blob/64916fa1fdbc3e444a652eca204d5669058ee321/src/problems/finals/12/14.md}{逐题来源及版本}。

\begin{solve}

把红球个数写成独立示性变量之和，均值给出主要规模、方差给出波动规模。令 $X_k$ 为第 $k$ 盒是否抽到红球的示性变量，则 $X_k\sim\operatorname{Bernoulli}(1/k)$ 且 $S_n=\sum_{k\leq n}X_k$。记
\[
\mu_n=ES_n=\log n+O(1),\qquad
s_n^2=\operatorname{Var}S_n=\sum_{k=1}^n\left(\frac1k-\frac1{k^2}\right)=\log n+O(1).
\]

（1）$E(S_n-\log n)^2=s_n^2+(\mu_n-\log n)^2=O(\log n)$，故 Chebyshev 不等式给出 $S_n/\log n\to1$ 依概率成立。

（2）因 $|X_k-EX_k|\leq1$ 且 $s_n\to\infty$，Lindeberg 条件中的截断项最终全为零，故 $(S_n-\mu_n)/s_n\Rightarrow N(0,1)$。再用 $s_n/\sqrt{\log n}\to1$ 及 $(\mu_n-\log n)/\sqrt{\log n}\to0$，由 Slutsky 定理得到结论。

（3）令 $Z_n=|S_n/\log n-1|^r/(1+|S_n/\log n-1|^r)$。由 （1）及连续映射定理，$Z_n\to0$ 依概率成立，且 $0\leq Z_n\leq1$。因此 $EZ_n\leq\varepsilon+P(Z_n>\varepsilon)\to\varepsilon$，令 $\varepsilon\downarrow0$ 即得所求。

\end{solve}

\begin{remark}[方法与条件]

$S_n$ 也与均匀随机排列的循环数同分布：依次加入数字 $k$ 时，在 $k$ 种等可能选择中，一种新建循环，其余把它接在已有某个数字后面；独立选择使新增循环数为 $\operatorname{Bernoulli}(1/k)$。本题因此同时描述随机排列的循环数，见\hyperref[app:probability-limit]{方法卷：概率极限定理}。

\end{remark}

\section{第13届决赛：12题}\label{pr:fin13-high:12}

\begin{example}[公开题解转录]\label{q:fin13-high:12}

记 $p$ 是整数，$p \geq 2, \beta_0\in\mathbb R^p$ 为固定参数，$X$ 是 $p$ 维随机向量，服从正态分布 $N(0, \Sigma)$。其中，$\Sigma$ 是对称 $p$ 阶正定矩阵。随机误差 $\varepsilon$ 是 1 维随机变量，服从正态分布 $N(0, \sigma^2)$，且与 $X$ 独立，其中 $\sigma^2 > 0$。记 $\operatorname{sgn}(\cdot)$ 为符号函数，

\[
Y = \operatorname{sgn}\left(X^{\mathrm{T}}\beta_0 + \varepsilon\right)
\]

是 1 维随机变量。定义

\[
\beta^* = \operatorname*{arg}\min_{\beta \in \mathbb{R}^p} E\left(Y - X^{\mathrm{T}}\beta\right)^2,
\]

其中 $\operatorname{arg}\min$ 是使得目标函数达到最小值的点。

（1）请给出 $\beta^*$ 的显式表示式。

（2）证明：$\operatorname{cov}(X, Y) = -2 \operatorname{cov}\left(X, F\left(-X^{\mathrm{T}}\beta_0\right)\right)$，其中 $F(\cdot)$ 表示 $\varepsilon$ 的分布函数。

（3）对任意的 $p$ 维参数向量 $\alpha, \beta$（其中 $\beta\ne0$），当 $X \sim N(0, \Sigma)$ 时，证明：

\[
E\left(\alpha^{\mathrm{T}}X \mid \beta^{\mathrm{T}}X\right) = \left(\alpha^{\mathrm{T}}\Sigma\beta\right)\left(\beta^{\mathrm{T}}\Sigma\beta\right)^{-1}\beta^{\mathrm{T}}X.
\]

（4）证明：$\beta^*$ 与 $\beta_0$ 成比例。要求找到比例因子的显式表达式。

\end{example}

\href{https://github.com/lzxcr/cmc-mathematics/blob/64916fa1fdbc3e444a652eca204d5669058ee321/src/problems/finals/13/15.md}{逐题来源及版本}。

\begin{solve}

（1）最小二乘是把 $Y$ 投影到 $X$ 的线性组合上。目标函数为 $\mathbb EY^2-2\beta^{\mathrm T}\mathbb E(XY)+\beta^{\mathrm T}\Sigma\beta$。由 $\Sigma>0$，唯一极小点为
\[
\beta^*=\Sigma^{-1}\mathbb E(XY)=\Sigma^{-1}\operatorname{cov}(X,Y).
\]
（2）独立性给出 $\mathbb E(Y\mid X)=1-2F(-X^{\mathrm T}\beta_0)$。乘以 $X$ 后取期望，利用 $\mathbb EX=0$，即得所需协方差等式。

（3）当 $\beta\ne0$ 时，令 $c=(\alpha^{\mathrm T}\Sigma\beta)/(\beta^{\mathrm T}\Sigma\beta)$。联合正态变量 $\alpha^{\mathrm T}X-c\beta^{\mathrm T}X$ 与 $\beta^{\mathrm T}X$ 不相关，故独立，且前者均值为零。这就证明条件期望公式。若 $\beta=0$，条件期望直接为零。

（4）令 $Z=X^{\mathrm T}\beta_0+\varepsilon$，$v=\beta_0^{\mathrm T}\Sigma\beta_0+\sigma^2>0$。同样的正态投影公式给出
\[
\mathbb E(X\mid Z)=\frac{\Sigma\beta_0}{v}Z.
\]
由 $Z\sim N(0,v)$，直接积分得 $\mathbb E|Z|=\sqrt{2v/\pi}$。因此
\[
\mathbb E(XY)=\frac{\Sigma\beta_0}{v}\mathbb E|Z|
=\sqrt{\frac2{\pi v}}\,\Sigma\beta_0,
\qquad
\boxed{\beta^*=\sqrt{\frac2{\pi(\beta_0^{\mathrm T}\Sigma\beta_0+\sigma^2)}}\,\beta_0}.
\]
此式也包括 $\beta_0=0$。

\end{solve}

\begin{remark}[方法与条件]

最小二乘只取决于 $\mathbb E(XY)$，正态回归把这个向量期望化为一维积分。更一般地，若 $Y=h(X^{\mathrm T}\beta_0,\varepsilon)\in L^2$，噪声独立于 $X$，且 $\beta_0\ne0$，则
\[
\beta^*=\frac{\mathbb E[(X^{\mathrm T}\beta_0)Y]}{\beta_0^{\mathrm T}\Sigma\beta_0}\,\beta_0.
\]
这里先给定指标 $X^{\mathrm T}\beta_0$，正交于它的正态残差与指标及噪声均独立。条件投影与正态回归的证明见\hyperref[app:probability-gaussian]{方法卷：正态条件结构}。

\end{remark}

\section{第14届决赛：12题}\label{pr:fin14-high:12}

\begin{example}[公开题解转录]\label{q:fin14-high:12}

设 $\left(T_n\right)_{n \in \mathbb{N}}$ 为一列独立同分布随机变量且其共同分布为参数为 1 的指数分布。定义 $S_0 = 0, S_n = T_1 + \cdots + T_n, n \geq 1$ 以及

\[
N_t := \max \left\{n \geq 0 : S_n \leq t\right\},\quad t \geq 0.
\]

（i）求 $S_n$ 的概率分布函数以及 $N_t$ 的概率分布列；

（ii）假定 $f \in L^1([0,1])$ 为一可积函数，试证明

\[
\mathbb{E}\left(\int_0^1 f(s) \mathrm{d}N_s\right) = \mathbb{E}\left(\sum_{n=1}^{\infty} f\left(S_n\right)1_{\{S_n \leq 1\}}\right) = \int_0^1 f(s) \mathrm{d}s;
\]

（iii）基于上述等式构造计算积分 $\int_0^1 f(s) \mathrm{d}s$ 的蒙特卡洛算法。

\end{example}

\href{https://github.com/lzxcr/cmc-mathematics/blob/64916fa1fdbc3e444a652eca204d5669058ee321/src/problems/finals/14/15.md}{逐题来源及版本}。

\begin{solve}

（i）等待时间相加对应密度卷积。由 $g_1(s)=e^{-s}1_{\{s>0\}}$，归纳得 $g_n(s)=s^{n-1}e^{-s}1_{\{s>0\}}/(n-1)!$。积分后，$t\ge0$ 时
\[
\mathbb P(S_n\le t)=1-e^{-t}\sum_{j=0}^{n-1}\frac{t^j}{j!},
\qquad
\mathbb P(N_t=k)=\mathbb P(S_k\le t<S_{k+1})=e^{-t}\frac{t^k}{k!},
\]
其中 $k=0,1,\dots$；$S_n$ 的分布函数在 $t<0$ 时为零。

（ii）期望到达点的密度是 $\sum_{n\ge1}g_n(s)=1$，这正是无偏性的来源。每次跳跃的大小为 $1$，故路径上积分等于题设随机和。先对 $|f|$ 用 Tonelli 定理：
\[
\mathbb E\sum_{n\ge1}|f(S_n)|1_{\{S_n\le1\}}
=\int_0^1|f(s)|e^{-s}\sum_{n\ge1}\frac{s^{n-1}}{(n-1)!}\,\mathrm ds
=\|f\|_1<\infty.
\]
因此可对 $f$ 合法交换期望、求和与积分，得到期望为 $\int_0^1f$。

（iii）独立生成指数等待时间并累加，至首次超过 $1$ 为止；对此前到达时刻求和，得到一次样本 $Z=\sum_{S_n\le1}f(S_n)$。重复 $m$ 次取均值 $m^{-1}\sum_{j=1}^mZ_j$，它无偏且由强大数定律几乎处处收敛到目标积分。

等待时间、Poisson 计数与随机积分是同一过程的三种描述；变换及矩公式见\hyperref[app:probability-poisson]{方法卷：Poisson 到达与随机和}。

\end{solve}

\section{第15届决赛：12题}\label{pr:fin15-high:12}

\begin{example}[公开题解转录]\label{q:fin15-high:12}

考虑二维整数格点 $\mathbb{Z}^{2} = \left\{k\mathrm{e}_{1} + j\mathrm{e}_{2} : k,j \in \mathbb{Z}\right\}$，其中 $\mathrm{e}_{1} = (0,1)$，$\mathrm{e}_{2} = (1,0)$。给定 $\varepsilon > 0$，设 $\left(\xi_{n}^{\varepsilon}\right)_{n \in \mathbb{N}}$ 为一列取值于 $\varepsilon\mathbb{Z}^{2}$ 上的独立同分布随机变量，且其共同分布为 $\mathbb{P}\left\{\xi_{n}^{\varepsilon} = \varepsilon\mathrm{e}_{i}\right\} = \mathbb{P}\left\{\xi_{n}^{\varepsilon} = -\varepsilon\mathrm{e}_{i}\right\} = \frac{1}{4}$（$i = 1,2$）。$\left(N_{t}\right)_{t \geq 0}$ 为一独立标准 Poisson 过程，即 $\mathbb{P}\left(N_{t} = k\right) = \frac{t^{k}\mathrm{e}^{-t}}{k!}$。定义 $S_{0}^{\varepsilon} = 0$ 以及 $S_{n}^{\varepsilon} := \xi_{1}^{\varepsilon} + \cdots + \xi_{n}^{\varepsilon}$（$n \in \mathbb{N}$），$Z_{t}^{\varepsilon} := S_{N_{t/\varepsilon^{2}}}^{\varepsilon}$（$t > 0$）。设 $\varphi: \mathbb{R}^{2} \rightarrow \mathbb{R}$ 为一有界连续函数，定义 $u_{\varepsilon}(t,z) := P_{t}^{\varepsilon}\varphi(z) := \mathbb{E}\varphi\left(Z_{t}^{\varepsilon} + z\right)$。
（1）计算 $Z_{t}^{\varepsilon}$ 的特征函数。
（2）证明：对任意 $t,s \geq 0$ 以及 $z \in \varepsilon\mathbb{Z}^{2}$，$P_{t+s}^{\varepsilon}\varphi(z) = P_{t}^{\varepsilon} \circ P_{s}^{\varepsilon}\varphi(z)$。
（3）验证 $\partial_{t}u_{\varepsilon} = \Delta_{\varepsilon}u_{\varepsilon}/4$，其中 $\Delta_{\varepsilon}f(z) := \sum_{i=1}^{2}\frac{f\left(z + \varepsilon\mathrm{e}_{i}\right) + f\left(z - \varepsilon\mathrm{e}_{i}\right) - 2f(z)}{\varepsilon^{2}}$。
（4）证明：对任意 $t > 0$ 以及 $z \in \mathbb{R}^2$，$\lim_{\varepsilon \downarrow 0}u_{\varepsilon}(t,z) = \mathbb{E}\varphi\left(B_{t/2} + z\right) = u(t,z)$，其中 $\left(B_{t}\right)_{t \geq 0}$ 为二维标准布朗运动。从而 $\partial_{t}u = \Delta u/4$，其中 $\Delta$ 为 Laplace 算子。

\end{example}

\href{https://github.com/lzxcr/cmc-mathematics/blob/64916fa1fdbc3e444a652eca204d5669058ee321/src/problems/finals/15/15.md}{逐题来源及版本}。

\begin{solve}

（1）一步跳跃的特征函数为 $\psi_\varepsilon(\theta)=[\cos(\varepsilon\theta_1)+\cos(\varepsilon\theta_2)]/2$。对 Poisson 跳跃次数条件化得
\[
\mathbb E e^{i\theta\cdot Z_t^\varepsilon}
=\exp\left\{\frac t{\varepsilon^2}
\left[\frac{\cos(\varepsilon\theta_1)+\cos(\varepsilon\theta_2)}2-1\right]\right\}.
\]
（2）复合 Poisson 过程具有平稳独立增量，故 $Z_{t+s}^\varepsilon$ 与 $Z_t^\varepsilon+\widetilde Z_s^\varepsilon$ 同分布，右边两项独立。先对第二个独立增量取期望，再对第一个取期望，即给出 $P_{t+s}^\varepsilon=P_t^\varepsilon P_s^\varepsilon$。

（3）令 $T_\varepsilon\varphi(z)=\frac14\sum_{i=1}^2[\varphi(z+\varepsilon e_i)+\varphi(z-\varepsilon e_i)]$。在有界函数空间上 $\|T_\varepsilon\|\le1$，而对跳数条件化得到 $P_t^\varepsilon=e^{-t/\varepsilon^2}\sum_{k\ge0}(t/\varepsilon^2)^kT_\varepsilon^k/k!$。该级数按算子范数局部一致收敛且可逐项求导，所以
\[
P_t^\varepsilon=e^{(t/\varepsilon^2)(T_\varepsilon-I)},\qquad
\partial_tP_t^\varepsilon\varphi=\frac{T_\varepsilon-I}{\varepsilon^2}P_t^\varepsilon\varphi
=\frac14\Delta_\varepsilon u_\varepsilon.
\]
（4）由 $\cos v=1-v^2/2+o(v^2)$，特征函数趋于 $e^{-t|\theta|^2/4}$，即 $B_{t/2}$ 的特征函数。Lévy 连续性定理及 $\varphi$ 有界连续给出所需期望收敛。

最后直接用极限的 Gaussian 积分表示
\[
u(t,z)=\frac1{\pi t}\int_{\mathbb R^2}e^{-|w-z|^2/t}\varphi(w)\,\mathrm dw
\]
可得 $t>0$ 时 $\partial_tu=\Delta u/4$。微分作用在热核上；在远离 $t=0$ 的紧时间区间内，其导数由可积的 Gaussian 函数乘多项式控制，因此只要求 $\varphi$ 有界，无须对它求导。

\end{solve}

\begin{remark}[方法与条件]

Poisson 化把一步转移算子变成连续时间半群；扩散尺度由单步协方差与单位时间跳数共同决定。特征函数和生成算子的关系见\hyperref[app:probability-poisson]{方法卷：Poisson 随机和}。

\end{remark}

\section{第16届决赛：12题}\label{pr:fin16-high:12}

\begin{example}[公开题解转录]\label{q:fin16-high:12}

设 $X_1,X_2,\cdots,X_n$ 是独立同分布的随机变量，且 $X_i$ 服从参数为 1 的指数分布。定义 $Y = \min\{X_1,X_2,\cdots,X_n\}$，$Z_i = X_i - Y$ 以及 $S_n = \sum_{i=1}^{n}X_i$。
（1）分别求 $Y,S_n$ 的概率分布密度以及事件 $\{Y = X_1\}$ 的概率。
（2）证明在 $Y = X_1$ 的条件下 $Z_2,\cdots,Z_n$ 独立且服从参数为 1 的指数分布。即对任意的 $x_i > 0$，$\mathbb{P}(Z_2 \leq x_2,\cdots,Z_n \leq x_n \mid Y = X_1) = \prod_{i=2}^{n}\mathbb{P}(X_i \leq x_i)$。
（3）令 $T := S_n - nY$，证明 $T$ 和 $S_{n-1}$ 同分布且 $T$ 与 $Y$ 独立。
（4）求 $\frac{Y}{S_n}$ 的概率分布密度。

\end{example}

\href{https://github.com/lzxcr/cmc-mathematics/blob/64916fa1fdbc3e444a652eca204d5669058ee321/src/problems/finals/16/15.md}{逐题来源及版本}。

\begin{solve}

（1）独立性给出 $\mathbb P(Y>y)=e^{-ny}$，故 $Y$ 的密度为 $ne^{-ny}1_{\{y>0\}}$。又 $S_n\sim\operatorname{Gamma}(n,1)$，密度为 $s^{n-1}e^{-s}1_{\{s>0\}}/(n-1)!$。最小值几乎处处唯一，由对称性，$\mathbb P(Y=X_1)=1/n$。

（2）令 $J$ 为取得最小值的指标。给定 $J=j$，换元 $X_j=y$、$X_i=y+z_i$（$i\ne j$），Jacobian 为 $1$。未归一化联合密度为
\[
e^{-ny}\prod_{i\ne j}e^{-z_i},\qquad y>0,\ z_i>0.
\]
除以 $\mathbb P(J=j)=1/n$ 后，它分解为 $ne^{-ny}\prod_{i\ne j}e^{-z_i}$。故条件下 $Y$ 与全部剩余等待时间相互独立，且各 $Z_i\sim\operatorname{Exp}(1)$。取 $j=1$ 即为所问。

（3）设 $n\ge2$。在每个 $J=j$ 条件下，$T=\sum_{i\ne j}Z_i\sim\operatorname{Gamma}(n-1,1)$，并与 $Y\sim\operatorname{Exp}(n)$ 独立。该联合分布不依赖 $j$，故无条件结论相同，且 $T\overset d=S_{n-1}$。

（4）对 $n\ge2$，$nY\sim\operatorname{Gamma}(1,1)$ 与 $T\sim\operatorname{Gamma}(n-1,1)$ 独立。令 $U=nY/(nY+T)$、$V=nY+T$，反变换为 $(nY,T)=(UV,(1-U)V)$，Jacobian 为 $V$。联合密度成为 $v^{n-1}e^{-v}(1-u)^{n-2}/(n-2)!$；对 $v>0$ 积分得 $p_U(u)=(n-1)(1-u)^{n-2}$（$0<u<1$），即 $U\sim\operatorname{Beta}(1,n-1)$。再按比例缩放，得到
\[
\boxed{p_{Y/S_n}(x)=n(n-1)(1-nx)^{n-2}1_{\{0<x<1/n\}}}.
\]
若 $n=1$，则 $Y/S_n\equiv1$，为点质量分布，没有 Lebesgue 密度；此时可约定 $S_0=T=0$。

\end{solve}

\begin{remark}[方法与条件]

指数分布的无记忆性在此表现为联合密度的因子分解；Gamma 和的比例又产生 Beta 分布。变量变换和条件分解的统一做法见\hyperref[app:probability-transforms]{方法卷：变量变换}及\hyperref[app:probability-convolution]{方法卷：卷积}。

\end{remark}

\section{第17届决赛：12题}\label{pr:fin17-high:12}

\begin{example}[公开题解转录]\label{q:fin17-high:12}

设 $X,Y$ 是离散随机变量，取值于有限集 $\mathcal X,\mathcal Y$。其概率分布为 $p(x)=P(X=x)$，$p(y)=P(Y=y)$，联合分布为 $p(x,y)=P(X=x,Y=y)$，条件分布为 $p(x|y)=P(X=x|Y=y)$（当 $p(y)>0$ 时）。熵、联合熵和条件熵定义如下：

\[
H(X):=-\sum_{x\in\mathcal X}p(x)\ln p(x),\qquad
H(X,Y):=-\sum_{x\in\mathcal X}\sum_{y\in\mathcal Y}p(x,y)\ln p(x,y),
\]

\[
H(X|Y):=-\sum_{x,y}p(x,y)\ln p(x|y),
\]

约定 $0\ln0=0$。证明：

（1）链式法则：$H(X,Y)=H(X)+H(Y|X)$；

（2）条件熵递减：$H(X|Y)\le H(X)$，且等号成立当且仅当 $X$ 与 $Y$ 独立。

\end{example}

\href{https://github.com/lzxcr/cmc-mathematics/blob/64916fa1fdbc3e444a652eca204d5669058ee321/src/problems/finals/17/15.md}{逐题来源及版本}。

\begin{solve}

（1）在 $p(x,y)>0$ 的位置，$p(x,y)=p(x)p(y\mid x)$。取对数，乘以 $-p(x,y)$ 后求和，得到
\[
H(X,Y)=H(X)+H(Y\mid X).
\]
零概率项按约定忽略，故无需给零概率条件事件定义条件分布。

（2）同样展开可得
\[
H(X)-H(X\mid Y)
=\sum_{p(x,y)>0}p(x,y)\log\frac{p(x,y)}{p(x)p(y)}
=D(p_{XY}\Vert p_Xp_Y).
\]
Gibbs 不等式说明右端非负：由 $-\log t\ge1-t$，对概率分布 $p,q$，
\[
\sum_{p_i>0}p_i\log\frac{p_i}{q_i}
\ge1-\sum_{p_i>0}q_i\ge0,
\]
且等号恰在两分布完全相同时成立。因此条件熵不增，等号当且仅当 $p(x,y)=p(x)p(y)$，也就是独立。

\end{solve}

\begin{remark}[方法与条件]

互信息同时等于条件化减少的熵，以及联合分布相对独立分布的偏离。非负性与等号条件来自同一个对数不等式；相关概念及条件互信息的推广见\hyperref[app:probability-entropy]{方法卷：熵与凸性}。

\end{remark}





